\chapter{How to cite these notes}
\label{app:citing}

These notes are revised: chapters gain material, the evidence is refreshed, and the version
number on the title page moves. A reference to them therefore has to answer the same question
that Chapter~\ref{ch:ardc} puts to any research artefact --- does it name \emph{the work}, which
will keep changing, or \emph{one exact version}, which will not? Both are legitimate. What is not
legitimate is a reference that leaves the reader unable to tell which was meant.

\begin{callout}{The short answer}
\small
Cite the HAL record \textbf{without} a version number:
\begin{center}
\url{https://hal.science/hal-05734828}
\end{center}
That address always resolves to the most recent version. A reader who follows it in five years
gets whatever these notes have become by then, which is what a reference to a living course note
should do.
\end{callout}

\section{Citing the work}

This is the form to use in a reading list, a syllabus, a policy document, or a paper that
recommends the notes as a whole. In plain text:

\begin{quote}\footnotesize
Roberto Di Cosmo. \emph{The Source Code of Science}. Doctoral. Pise (Italia), Italy.
2026. \url{https://hal.science/hal-05734828}
\end{quote}

\noindent and in Bib\TeX{}, adapted from the entry HAL generates, with the version-specific fields
left out on purpose:

\begin{quote}\footnotesize
\begin{verbatim}
@unpublished{dicosmo:hal-05734828,
  TITLE   = {{The Source Code of Science}},
  AUTHOR  = {Di Cosmo, Roberto},
  URL     = {https://hal.science/hal-05734828},
  TYPE    = {Doctoral},
  ADDRESS = {Pise (Italia), Italy},
  YEAR    = {2026},
  HAL_ID  = {hal-05734828},
}
\end{verbatim}
\end{quote}

\noindent The \texttt{URL} carries no \texttt{v} suffix, and that is the whole point: HAL keeps
every version, and the unsuffixed identifier is the one that follows the work forward.

\section{Citing one exact version}

This is the form to use when the reference supports a claim --- when you quote a page, rely on a
figure, or say that these notes recommend some specific thing. A later version may say something
else, and a reader checking your claim must land on the text you actually read.

Three identifiers name an exact version, at three different levels:

\begin{description}[leftmargin=1.6em,itemsep=2pt]
\item[The HAL record, versioned.] Append the version suffix HAL assigns, as in
  \url{https://hal.science/hal-05734828v1}. Each deposit is a fixed document with a fixed cover
  page. HAL's own citation block for that version carries the page count and the version note,
  so take it from the record rather than retyping it.
\item[The release.] Every version is tagged in the source repository and published with its
  \textsc{pdf} attached:
  \url{https://github.com/rdicosmo/source-code-of-science/releases}. The tag (\texttt{v1.4},
  \texttt{v1.5}, \dots) is the human-readable name of the version.
\item[The source tree.] The \LaTeX{} sources, figures and illustrations of each released version
  are archived in Software Heritage and named by a qualified \textsc{swhid}. This is the
  identifier to use when what you need is the material these notes are built from, rather than a
  rendering of them. Version~1.5, for example, is named by
  \begin{quote}\footnotesize
  \swhreflong{swh:1:dir:769f7afe6cc52fe86bd63d35cb68e7513874e2fd;origin=https://github.com/rdicosmo/source-code-of-science;visit=swh:1:snp:b8715018e8f7ffab19b3f8a88fcd036fd2020ca8;anchor=swh:1:rev:7f22576697b3ee89eb96ac11839f724fe7db91b0}{swh:1:dir:769f7afe6cc52fe86bd63d35cb68e7513874e2fd;origin=https://github.com/rdicosmo/source-code-of-science;visit=swh:1:snp:b8715018e8f7ffab19b3f8a88fcd036fd2020ca8;anchor=swh:1:rev:7f22576697b3ee89eb96ac11839f724fe7db91b0}
  \end{quote}
  \vspace{-4pt}
  The part before the first semicolon, \texttt{swh:1:dir:769f7afe\dots}, is the identifier
  proper: it is computed from the content, so it can be recomputed and verified offline by the
  method of Chapter~\ref{ch:identifiers}. What follows it is context that helps a reader find the
  object again: where it came from, which visit saw it, and which revision it hangs from. Each
  released version has its own such identifier, recorded in that version's HAL record under
  related resources and in the release notes.
\end{description}

\noindent In Bib\TeX{}, the versioned form differs from the one above in two fields:

\begin{quote}\footnotesize
\begin{verbatim}
  URL         = {https://hal.science/hal-05734828v1},
  HAL_VERSION = {v1},
\end{verbatim}
\end{quote}

\section{Which document is which}

One caution, and it is the note's own lesson turned on itself. The \textsc{pdf} served by HAL is
\emph{not} byte-identical to the \textsc{pdf} attached to the corresponding release: HAL prepends
a cover page, which makes it a different digital object even though it is the same document to a
reader. A byte-stable copy of each release, exactly as built, is kept at
\url{https://dicosmo.org/Articles/2026-source-code-of-science.pdf}.

The \textsc{swhid} recorded with each version names the \emph{source tree}, not any of those
\textsc{pdf}s. This is not an oversight. A rendered document has no canonical byte form --- it can
be rebuilt, converted, or reprinted with a cover page, and remain the same work --- whereas its
source does, which is precisely why intrinsic identifiers fit code and fit sources, and sit
awkwardly on a printed artefact.

The example above names version~1.5. The version you are holding names itself as well:
Appendix~\ref{app:self} prints the identifier of the sources this copy was built from, and
explains how the build derives it. The distinction matters. A document cannot carry an intrinsic
identifier of \emph{itself}, since writing the identifier in changes the bytes it names. It can
carry the identifier of the \emph{sources it was built from}, which is a different object, and
the one you want in a citation anyway.

\takeaway{Cite the unsuffixed identifier for the work, the versioned one for a claim. A reference
that does not say which it means is not doing its job.}
